% The spoken body of the homily for the Twenty-fifth Sunday in Ordinary Time,
% Year A. Continuous preaching, ready to read aloud: no heading, no direction
% to the preacher, no citation inside the speech. The loci, the relation to
% the reviewed interpretations, the audience, the word count and the pace
% estimate stand in the terminal note, which no one reads out. The blank lines
% are the paragraph breaks a speaker recovers his place by, and the five
% vertical spaces divide the speech into its six movements: the line at
% evening and the complaint; the householder's answer, with the psalm's two
% words and the evil eye; the prophet's reason; what the wage is, and Paul;
% the two kinds of hearer; and the prayers of the Mass, with the second line.
%
% Scripture is spoken in the Douay--Rheims, the leaf's study translation, and
% the speech names that Bible once, just before the first verse is quoted. No oration or
% antiphon of the Missal is quoted: each is described.

Evening has come in the vineyard, and the labourers are standing in line for
their pay. At the back are the men who were hired at dawn. They have worked
twelve hours, and the heat of the day is still in their shoulders. At the
front are the men who came an hour ago, when the light was already going.
The steward begins with them, and into each of their hands he puts a
denarius, the coin agreed that morning as the wage for a whole day. The old
Douay Bible calls it a penny.

You can feel the arithmetic begin at the back of the line. If one hour is
worth a denarius, what must twelve be worth? Then their own turn comes, and
each of them receives a denarius. And they say what any of us would say. In
that Bible's words: ``These last have worked but one hour, and thou hast made
them equal to us, that have borne the burden of the day and~the~heats.''

If we are honest, we are standing with them. Their complaint sounds like
plain fairness. So ask the question the parable asks. What, exactly, have
these men lost?

Nothing. They agreed for a denarius, and a denarius is in their hand. What
pains them is nothing that was done to them. It is the good that was done to
somebody else.

\bigskip

That is where the householder's answer goes. ``Friend, I do thee no wrong:
didst thou not agree with me for a penny? Take what is thine, and go thy way:
I will also give to this last even as to thee.'' Saint Augustine, preaching
on this Gospel, supplies the words the householder leaves unspoken: ``I have
paid thee what I agreed for. To this other it is my will not to render a
payment, but to bestow a gift.'' Justice to the one, a gift to the other, and
no one robbed. Today's psalm says both things of the one Lord: he is
``gracious and merciful'', and he is ``just in all his ways''.

Then comes the last question: ``Is thy eye evil, because I am good?'' Saint
Thomas Aquinas says what an evil eye is. A man is properly called evil, he
says, when he grieves at goodness. And Saint Jerome finds the wish hidden
under the complaint: you do not so much want more for yourself as you want
the other man to receive nothing.

\bigskip

The first reading was chosen to stand beside this Gospel, and it gives the
householder's reason before the parable is told: ``My thoughts are not your
thoughts''. We reach for that line when God's ways bewilder us. But listen to
where the prophet puts it. ``Let the wicked forsake his way, and the unjust
man his thoughts, and let him return to the Lord, and he will have mercy on
him, and to our God: for he is bountiful to forgive. For my thoughts are not
your thoughts: nor your ways my ways, saith the Lord.'' The line is about
pardon. Saint Thomas gives the whole contrast in a breath: you think of
requital, of paying back; I think of mercy. The men at the back of the line
reasoned by the thoughts of men: more hours, more pay. The householder's
reckoning is as far above theirs as the heavens are above the earth.

\bigskip

But we have not yet asked what the coin is. Saint Augustine says it plainly:
that denarius is eternal life. The saints will shine, he says, some more and
some less; but in living for ever no one will live more than another. Saint
Jerome noticed what is stamped on a denarius: the figure of the king. You
have received the wage I promised you, he hears the householder say, my own
image and likeness; what more do you seek? And Augustine, preaching on
today's psalm, comes to its promise that the Lord is nigh ``to all that call
upon him in truth.'' To call on him in truth, he says, is to want the giver
more than his gifts. How much more blessed will you be, he asks, when he
gives you himself?

That is why this wage cannot be cut into shares. What God holds out to the
labourer of the eleventh hour is himself, and he has no smaller self to give
the latecomer. The man who thinks he has been paid short has not yet looked
at what is in his hand.

Saint Paul has looked. In the second reading he is the labourer those
grumblers might have been. Few have borne more of the day's burden, and he
writes as a prisoner, in chains. His wage is within reach: ``to die is
gain''; to be with Christ is ``a thing by far the better.'' And he is willing
to stay at his work: ``to abide still in the flesh is needful for you.'' The
grumblers reckon their hours against the latecomers. Paul adds hours for
their sake. He can, because for him the wage was never a sum: ``To me, to
live is Christ''. Whoever has Christ already loses nothing when Christ is
generous to someone else.

\bigskip

Saint John Chrysostom says that this parable is spoken to two kinds of
people. It is spoken to those who came to God early, ``that they may not be
proud, neither reproach those called at the eleventh hour''. And it is spoken
to those who come late, ``that they may learn that it is possible even in a
short time to recover all.'' Both kinds are here today.

Many of us were hired at dawn. Baptized as infants, we have kept the faith
through long years, and it has cost us something. Thank God for every hour
of it. But long service has a temptation of its own, and it is the evil eye.
It wakes when a man who mocked the Church for forty years is reconciled on
his deathbed; when someone who has carried none of our burdens is given
everything we have been given. Something in us begins to keep accounts.
Chrysostom is certain that there is no such grumbling among the saints:
heaven, he says, ``is pure from envy and jealousy.'' So the murmur in this
parable is no report from heaven. It is set before us now, while we can
still be healed of it.

And some part of us may still be standing in the market-place. We are here,
and yet there is a confession put off for years, a wrong never set right, a
call from God that we are keeping for later. Do not reason that the eleventh
hour pays as well as the first. Augustine saw that reasoning coming. The
householder has promised you a denarius if you come at the eleventh hour, he
said; but whether you will live even to the seventh, no one has promised
you. ``Why then dost thou put off him that calleth thee, certain as thou art
of the reward, but uncertain of the day?''

\bigskip

Look, last, at how this Mass prays. The antiphon the Missal sets at its
entrance is spoken in God's own voice. Before anyone has asked him for
anything, he names himself the salvation of his people, and promises to hear
them from whatever trouble they cry out of. The psalm takes him at his word:
``The Lord is nigh unto all them that call upon him''. All of them: the one
who has called for seventy years, and the one who calls today for the first
time. Then, in the opening prayer, we spoke to the God who has made his
whole law rest on the love of himself and of our neighbour. We asked that
those who keep his commandments may come to eternal life. We did not present
a bill. We asked. The keeping and the arriving both stood inside a prayer.

Before this Mass ends another line will form, here in this church. In it will
stand those who have borne the burden of the day and the heats, and those
who came in at the eleventh hour; and to each will be given the same. Not a
coin this time. The Lord himself, the pledge of that one wage. And when we
have received him, the Church will ask that what we hold in these mysteries
we may hold also in the way we live.

So carry one thing out of the door with you. This week someone will be given
a good that looks to you unearned: a mercy, a welcome, a second chance. When
it happens, hear the householder say to you, gently, ``Friend, I do thee no
wrong''. Then do what the grumblers could not do. Be glad of it. And look at
what is in your own hand. It bears the image of the King. What more do we
seek?
